\documentclass{article}
\usepackage[T1]{fontenc}
\usepackage{lmodern}
\usepackage{graphicx}
\usepackage{amsmath,amssymb}
\usepackage[a4paper,margin=2cm]{geometry}
\usepackage[absolute,overlay]{textpos}
\setlength{\TPHorizModule}{1mm}
\setlength{\TPVertModule}{1mm}
\usepackage[indonesian]{babel}
\usepackage{hyperref}
\setlength{\emergencystretch}{2em}
% unit: Halaman 46

\begin{document}

% === Halaman 46 (python/native) ===

Program enkripsi-dekripsi AES dengan Python

Input: Teks dari keyboard

46

from Crypto.Cipher import AES

from Crypto.Random import get\_random\_bytes

from Crypto.Util.Padding import pad, unpad

import base64

\# Panjang blok AES

BLOCK\_SIZE = 16

def enkripsi\_aes(plainteks, kunci):                     

    kunci = kunci.encode()                              \# Kodekan kunci dari string menjadi byte

    iv = get\_random\_bytes(BLOCK\_SIZE)                   \# Bangkitkan IV

    AEScipher = AES.new(kunci, AES.MODE\_CBC, iv)        \# Defenisikan AES mode CBC


    plainteks = pad(plainteks.encode(),BLOCK\_SIZE)      \# Lakukan padding pada plainteks      


    cipherteks = AEScipher.encrypt(plainteks)           \# Lakukan enkripsi plainteks 

    cipherteks = base64.b64encode(iv+cipherteks)        \# Kodekan iv+cipherteks dari byte ke kode base64

    return cipherteks

\end{document}
